\documentclass[10pt,a4paper]{amsart}
\usepackage[margin=19mm]{geometry}
\usepackage{amsmath,amssymb,mathtools}
\usepackage[scr=rsfs,cal=euler]{mathalfa}
\usepackage{xfrac}
\usepackage[unicode,hidelinks,bookmarks=false]{hyperref}
\newcommand{\ie}{i.e.\ }
\newcommand{\cf}{cf.\ }
\newcommand{\resp}{respectively\ }
\newcommand{\Subsection}{\subsection}
\providecommand{\nobreakdash}{\nobreak\hbox{-}}
\setlength{\emergencystretch}{2em}
\multlinegap0pt
\usepackage{bm}
\usepackage{bbm}
\usepackage{dsfont}
\usepackage{xspace}
\usepackage{cancel}
\usepackage{mathtools}
\usepackage{amsthm}
\makeatletter
\@ifundefined{subsection}{\newtheorem{subsection}{Subsection}}{}
\@ifundefined{lemma}{\newtheorem{lemma}[subsection]{Lemma}}{}
\@ifundefined{prop}{\newtheorem{prop}[subsection]{Proposition}}{}
\@ifundefined{theorem}{\newtheorem{theorem}[subsection]{Theorem}}{}
\makeatother
\newcommand{\D}[1]{\mathbb{#1}}% Doubled - blackboard bold - only caps, uas as \D{A}
\newcommand{\UU}{\bigcup}% big union
\newcommand{\gb}{\beta}
\newcommand{\ii}{\cap}% intersection
\newcommand{\set}[1]{\{#1\}}% set
\newcommand{\te}{\text}% same as \mathrm command.
\renewcommand{\gg}{\gamma}% old use >>
\title[Conditional Optimal Sets and the Quantization Coefficients f]{Conditional Optimal Sets and the Quantization Coefficients for Some Uniform Distributions}
\author{Evans Nyanney, Megha Pandey,, Mrinal Kanti Roychowdhury}
\date{Source: arXiv:2402.16783v3 (CC BY 4.0)}
\begin{document}
\maketitle

\noindent\emph{Source and licence.} Conditional Optimal Sets and the Quantization Coefficients for Some Uniform Distributions. Version arXiv:2402.16783v3 (CC BY 4.0), distributed under the Creative Commons Attribution 4.0 International licence, as recorded in the arXiv licence metadata for this exact version and shown on the abstract page https://arxiv.org/abs/2402.16783v3. Full rights notice: licenses/arxiv-2402.16783-CC-BY-4.0.txt.\par\smallskip

\noindent\textit{Editorial scope.} Counted unit: the Prop whose source label is \texttt{prop45} in the article (source0.tex lines 598--612, source label \texttt{prop45}), delivered together with the article's own complete original proof of it (source0.tex lines 614--629, one \texttt{proof} environment of 16 lines). No mathematics is written, repaired or re-derived: statement and proof are the article's own text line for line. Required context retained uncounted: theoM0 (lines 356--368); lemma34 (lines 588--590). Every definition, display and label of those lines is retained unchanged. Omitted from this package and not redistributed: 1--355, 369--587, 591--597, 613--613, 630--897. Nothing inside a retained excerpt is omitted or altered: every retained body is delivered exactly as the article's lines are written, so no omission receipt of any kind accompanies this package. Citations: the retained excerpts carry no citation command at all, so no bibliography entry of the article is transcribed and no reference is redistributed. Reference adaptation: none was needed and none was made; every reference inside the delivered unit resolves inside it, so no unresolved reference or citation is produced. Printed slips kept literal: no printed word, symbol, hypothesis or number of the counted statement or of its proof was altered. Layout and class: the article's own class is \texttt{amsart} with packages including amssymb, latexsym, amsmath, amscd, amsthm, graphicx, color, xy, pgf, tikz, mathrsfs, geometry, hyperref; the wrapper is the lane's pinned \texttt{amsart} editorial wrapper at 10pt on A4, which restates the theorem-like environments the retained text needs and adds the article's own macro definitions that the retained text needs (\texttt{\textbackslash D}, \texttt{\textbackslash UU}, \texttt{\textbackslash gb}, \texttt{\textbackslash gg}, \texttt{\textbackslash ii}, \texttt{\textbackslash set}, \texttt{\textbackslash te}), with extra wrapper packages (bm, bbm, dsfont, xspace, cancel, mathtools). The delivered numbering is the unit's own. Assets: no retained span contains a figure environment, a graphics-inclusion command, a PDF-page inclusion, a table or a TikZ picture; no image file is copied into this package, the package declares no asset dependency and no image file is redistributed with it. Licence: version arXiv:2402.16783v3 of this article is distributed under the Creative Commons Attribution 4.0 International licence, as recorded in the arXiv licence metadata for this exact version and shown on the abstract page https://arxiv.org/abs/2402.16783v3. A full rights notice is delivered at licenses/arxiv-2402.16783-CC-BY-4.0.txt. Characters: every retained excerpt is pure ASCII, so the delivered engine, XeLaTeX with its default Latin Modern text font, reports no missing character.
\begin{theorem}\label{theoM0}
Let $P$ be the uniform distribution on the line segment joining $a$ and $b$, where $a, b\in \D R$ with $a<b$. Then we have:  

$(i)$ the conditional optimal set of $n$-points with respect to the conditional set $\gb:=\set{a, b}$ is  
\[\Big\{a+\frac{j-1}{\ell-1}(b-a) : 1\leq j\leq n\Big\} \te{ with conditional quantization error } V_n=\frac {(b-a)^2}{12(n-1)^2};\]


$(ii)$ the conditional optimal set of $n$-points with respect to the conditional set $\gb:=\set{a}$ is  
\[\Big\{a+\frac{2(j-1)(b-a)}{2n-1}   : 1\leq j\leq n\Big\} \te{ with conditional quantization error } V_n=\frac {(b-a)^2}{3(2n-1)^2};\]

$(iii)$ the conditional optimal set of $n$-points with respect to the conditional set $\gb:=\set{b}$ is  
\[\Big\{a+\frac{(2j-1)(b-a)}{2n-1}   : 1\leq j\leq n\Big\} \te{ with conditional quantization error } V_n=\frac {(b-a)^2}{3(2n-1)^2}.\]
\end{theorem}

\begin{lemma} \label{lemma34} 
Let $\gg_n$ be a conditional optimal set of $n$-points for the uniform distribution $Q$ with respect to the conditional set $\set{c_j : 1\leq j\leq m+1}$ for any $n\geq m+1$. Let $n_j=\te{card}(\gg_n\ii [c_j, c_{j+1}])$ for $1\leq j\leq m$. Then, $n_j\geq 2$ and $ n_1+n_2+\cdots +n_m=n+m-1$, and $|n_i-n_{j}|=0 \te{ or } 1$  for all $1\leq i\neq  j\leq m$. 
\end{lemma} 

\begin{prop} \label{prop45} 
Let $\gg_n$ be a conditional optimal set of $n$-points and $V_n(Q)$ be the $n$th conditional quantization error for the uniform distribution $Q$ with respect to the conditional set $\set{c_j : 1\leq j\leq m+1}$ for any $n\geq m+1$. Let $n=m k +1+q$, where $0\leq q< m$. Then, if $q=0$,  
we have 
\begin{align} \label{Megha101} 
\gg_n& = \UU_{j=1}^{m}\{c_j+\frac{2(i-1)r\sin \frac{\pi}m}{k}  : 1\leq i\leq k+1\} \te{ with } V_n(Q)=\frac {r^2 \sin^2\frac{\pi} m}{3k^2}.
\end{align} 
On the other hand, if $0<q<m$, then there are $^mC_q$ possible sets $\gg_n$, one such set is given by 
\begin{equation}\label{Megha102} 
\begin{aligned} 
\gg_n  &= \Big(\UU_{j=1}^{q}\{c_j+\frac{2(i-1)r\sin \frac{\pi}m}{k+1} : 1\leq i\leq k+2\}\Big)\\
&\qquad \UU \Big(\UU_{j=q+1}^{m}\{c_j+\frac{2(i-1)r\sin \frac{\pi}m}{k}  : 1\leq i\leq k+1\}\Big)\\
\te{ with }    V_n(Q)&=\frac {r^2q\sin^2\frac{\pi} m}{3m(k+1)^2}+ \frac {r^2(m-q)\sin^2\frac{\pi} m}{3mk^2}.
 \end{aligned}
\end{equation}  
 \end{prop} 

\begin{proof} 
First assume that $q=0$, then we have $n=mk+1$, i.e., $\gg_n$ contains $k-1$ elements from each of the intervals $[c_j, c_{j+1}]$ except the boundary elements $c_j$ and $c_{j+1}$. Hence, if $n_j=\te{card}(\gg_n\ii [c_j, c_{j+1}])$, we have $n_j=k-1+2=k+1$. Hence, by $(i)$ of Theorem~\ref{theoM0}, we have 
\begin{align*}
\gg_n& = \UU_{j=1}^{m}\{c_j+\frac{i-1}{k}(c_{j+1}-c_{j}) : 1\leq i\leq k+1\}=\UU_{j=1}^{m}\{c_j+\frac{2(i-1)r\sin \frac{\pi}m}{k}  : 1\leq i\leq k+1\}
\end{align*}
\[ \te{ with } V_n=\sum_{j=1}^m \frac {(c_{j+1}-c_j)^2}{12mk^2}=\sum_{j=1}^m \frac {r^2 \sin^2\frac{\pi} m}{3mk^2}=\frac {r^2 \sin^2\frac{\pi} m}{3k^2}.\]
On the other hand, if $0<q<m$, then due to Lemma~\ref{lemma34}, we can assume that $\gg_n$ contains $k$ elements from each of the first $q$ intervals except the boundary elements, and $\gg_n$ contains $(k-1)$ elements from each of the remaining $m-q$ intervals except the boundary elements implying $n_1=n_2=\cdots=n_q=k+2$ and $n_{q+1}=n_{q+2}=\cdots=n_m=k+1$. Hence, the expressions for $\gg_n$ and the corresponding $n$th conditional quantization errors  are  obtained by  $(i)$ of Theorem~\ref{theoM0} as 
\begin{align*}
\gg_n  &= \Big(\UU_{j=1}^{q}\{c_j+\frac{2(i-1)r\sin \frac{\pi}m}{k+1} : 1\leq i\leq k+2\}\Big)\\
&\qquad \UU \Big(\UU_{j=q+1}^{m}\{c_j+\frac{2(i-1)r\sin \frac{\pi}m}{k}  : 1\leq i\leq k+1\}\Big)
\end{align*} 
 with  
 \[V_n=\sum_{j=1}^q \frac {(c_{j+1}-c_j)^2}{12 m(k+1)^2}+\sum_{j=q+1}^m \frac {(c_{j+1}-c_j)^2}{12mk^2}=\sum_{j=1}^q \frac {r^2\sin^2\frac{\pi} m}{3m(k+1)^2}+\sum_{j=q+1}^m \frac {r^2\sin^2\frac{\pi} m}{3mk^2}\] 
 yielding \[V_n=\frac {r^2q\sin^2\frac{\pi} m}{3m(k+1)^2}+ \frac {r^2(m-q)\sin^2\frac{\pi} m}{3mk^2}.\]
 Notice that if $n=mk+q$, the optimal set $\gg_n$ can be constructed in $^mC_q$ ways. 
\end{proof} 

\end{document}
